% Part: sets-functions-relations
% Chapter: size-of-sets-complete

\documentclass[../../../include/open-logic-chapter]{subfiles}

\begin{document}

\olchapter{sfr}{siz}{সেটের আকার}

\begin{editorial}
এই অধ্যায়ে তালিকায়ন, গণনীয়তা ও অগণনীয়তা আলোচনা করা হয়েছে। কয়েকটি অংশের দুটি করে সংস্করণ আছে: একটি অপেক্ষাকৃত প্রাথমিক সংস্করণে তালিকায়নকে তালিকা, অথবা $\PosInt$ থেকে সমাপতিত অপেক্ষক হিসেবে ধরা হয়েছে; আর অপেক্ষাকৃত বিমূর্ত সংস্করণে তালিকায়নকে $\Nat$-এর সঙ্গে একৈক সমাপতিত অপেক্ষক হিসেবে সংজ্ঞায়িত করা হয়েছে।
\end{editorial}

\olimport{introduction}

\olimport{enumerability}

\olimport{zig-zag}

\olimport{pairing}

\olimport{pairing-alt}

\olimport{non-enumerability}

\olimport{reduction}

\olimport{equinumerous-sets}

\olimport{comparing-size}

\olimport{schroder-bernstein}

\begin{editorial}
পরবর্তী \olref[sfr][siz][enm-alt]{sec},
\olref[sfr][siz][nen-alt]{sec} ও \olref[sfr][siz][red-alt]{sec} হলো যথাক্রমে \olref[sfr][siz][enm]{sec},
\olref[sfr][siz][nen]{sec} ও \olref[sfr][siz][red]{sec}-এর বিকল্প সংস্করণ। টিম বাটন তাঁর \emph{Open Set Theory} পাঠ্যে ব্যবহারের জন্য এগুলি রচনা করেছেন। এগুলি কিছুটা বেশি অগ্রসর এবং সেটতত্ত্বের প্রসঙ্গে অধিক উপযোগী তালিকায়নযোগ্যতার একটি ভিন্ন সংজ্ঞা ব্যবহার করে (অর্থাৎ তালিকায় প্রকাশযোগ্য হওয়া বা $\PosInt$ থেকে কোনো সমাপতিত অপেক্ষকের বিস্তৃতি হওয়ার পরিবর্তে $\Nat$ বা তার কোনো প্রারম্ভিক খণ্ডের সঙ্গে একৈক সমাপতিত অপেক্ষক থাকা)।
\end{editorial}

\olimport{enumerability-alt}
\olimport{non-enumerability-alt}
\olimport{reduction-alt}

\OLEndChapterHook

\end{document}
